\documentclass[12pt,fullpage]{article}

\usepackage{amsthm}
\usepackage{amssymb}
\usepackage[fleqn]{amsmath}
\usepackage{enumitem}

\newtheorem{theorem}{Theorem}[section]
\newtheorem{lemma}[theorem]{Lemma}
\newtheorem{proposition}[theorem]{Proposition}
\newtheorem{corollary}[theorem]{Corollary}

\theoremstyle{definition}
\newtheorem{definition}[theorem]{Definition}
\newtheorem{remark}[theorem]{Remark}

\newcommand{\A}{\mathcal{A}}
\newcommand{\B}{\mathcal{B}}
\newcommand{\N}{\mathbb{N}}
\newcommand{\Z}{{\mathbb{Z}}}
\newcommand{\AUT}{\operatorname{AUT}}
\renewcommand{\phi}{\varphi}

\begin{document}

\title{Toward Question~1.2: a verified reduction, two refuted candidates,
and the precise obstruction}
\author{solver-agent}
\maketitle

\begin{abstract}
We work on Question~1.2 of Kalociński--Slaman: \emph{is there a computably
represented, computably $\AUT$-countable-on-a-cone structure $\A$ such that for
every finite parameter tuple $\bar a$ there is an automorphism $\pi$ that is not
$\Sigma^{in}_1$-definable from $\bar a\cup\pi(\bar a)$?}
We isolate and prove a robust \emph{embedding criterion}
(Lemma~\ref{lem:embcrit}) that converts the negative clause into a statement
about self-embeddings, and we derive the necessary conditions any witnessing
structure must satisfy (Corollary~\ref{cor:spec}).  We then analyse the two most
natural candidate constructions --- a directed/anchored spine
(``$\AUT\cong\Z$'') and an undirected coupling (``$\AUT\cong D_\infty$'') --- and
\emph{prove that both fail} the negative clause, by exhibiting explicit
existential definitions that defeat it.  These refutations locate the exact
obstruction, which we abstract in Lemma~\ref{lem:obstruction}: any
\emph{position-aligning} coupling strong enough to force a countable automorphism
group automatically leaks the orientation/value datum to a $\Sigma^{in}_1$
family.  We then formulate precisely the structural target a positive solution
must hit (Proposition~\ref{prop:target}).
\textbf{We do not claim a complete positive or negative answer to
Question~1.2; the construction realizing the target is left open and the issue is
reported as \emph{in progress}.}  All computations of automorphism groups, and
both refutations, are complete proofs and have been cross-checked numerically on
finite truncations.
\end{abstract}

\section{Conventions and the two clauses}

Throughout, ``$\Sigma^{in}_1$-definable from a finite tuple $\bar c$'' has the
meaning of the problem statement: $f\colon\A\to\A$ is
$\Sigma^{in}_1$-definable from $\bar c$ if there is a countable family of
existential formulas $(\phi_i)_{i\in\N}$ with
\begin{equation}\tag{$\ast$}\label{star}
f(x)=y\quad\Longleftrightarrow\quad \exists i\;\bigl(\A\models\phi_i(\bar c,x,y)\bigr)
\qquad(\forall x,y\in\A).
\end{equation}
By Theorem~2.3 of the problem, $\A$ is computably $\AUT$-countable on a cone iff
\begin{itemize}
\item[\textbf{(P1)}] every $\pi\in\AUT(\A)$ is $\Sigma^{in}_1$-definable from
  \emph{some} finite tuple $\bar c_\pi$.
\end{itemize}
The remaining requirement of Question~1.2 is the negative clause
\begin{itemize}
\item[\textbf{(P2)}] for every finite tuple $\bar a$ there is
  $\pi\in\AUT(\A)$ that is \emph{not} $\Sigma^{in}_1$-definable from
  $\bar a\cup\pi(\bar a)$.
\end{itemize}
Question~1.2 asks for a computably represented $\A$ with countable
automorphism group satisfying both (P1) and (P2).  Countability of $\AUT(\A)$
is required so that Theorems~2.2/2.3 apply; we verify it for each candidate.

\section{The embedding criterion (verified)}\label{sec:crit}

The decisive tool is that existential formulas are preserved by embeddings.  A
\emph{self-embedding} of $\A$ is an injective map $h\colon\A\to\A$ that
preserves \emph{and reflects} every relation and commutes with the functions and
constants; equivalently, $h$ is an isomorphism of $\A$ onto an induced
substructure of $\A$.  Automorphisms are exactly the surjective self-embeddings.

\begin{lemma}[Embedding criterion for non-definability]\label{lem:embcrit}
Let $\bar c$ be a finite tuple in $\A$ and $\pi\colon\A\to\A$ any function.
Suppose there is a self-embedding $h\colon\A\to\A$ with $h(c)=c$ for every entry
$c$ of $\bar c$, and a point $x_0\in\A$ with
$h\bigl(\pi(x_0)\bigr)\neq\pi\bigl(h(x_0)\bigr)$.
Then $\pi$ is \emph{not} $\Sigma^{in}_1$-definable from $\bar c$.
\end{lemma}

\begin{proof}
Suppose, for contradiction, that $(\phi_i)_{i\in\N}$ witnesses \eqref{star} for
$\pi$ over $\bar c$.  Put $y_0:=\pi(x_0)$.  By \eqref{star} some $\phi_i$
satisfies $\A\models\phi_i(\bar c,x_0,y_0)$.  As $\phi_i$ is existential and $h$
is a self-embedding, existential formulas are preserved under $h$, so
$\A\models\phi_i\bigl(h(\bar c),h(x_0),h(y_0)\bigr)$.  Since $h$ fixes $\bar c$
pointwise, $h(\bar c)=\bar c$, giving $\A\models\phi_i(\bar c,h(x_0),h(y_0))$.
By the ``$\Leftarrow$'' direction of \eqref{star} this forces
$\pi\bigl(h(x_0)\bigr)=h(y_0)=h\bigl(\pi(x_0)\bigr)$, contradicting the
hypothesis.
\end{proof}

\begin{remark}[Why automorphisms alone do not suffice]\label{rem:autos}
Taking $h$ to be an \emph{automorphism} requires $h\in\AUT(\A)$ fixing $\bar c$
pointwise yet not commuting with $\pi$.  In both candidates below the pointwise
stabilizer of a generic finite tuple is trivial (the action is essentially
free), so automorphisms give nothing; the force of Lemma~\ref{lem:embcrit} lies
entirely in admitting \emph{proper} (non-surjective) self-embeddings.  This is
the feature a positive construction must supply, and the feature the natural
candidates lack (\S\ref{sec:obstruction}).
\end{remark}

\begin{corollary}[Specification for a positive answer]\label{cor:spec}
To answer Question~1.2 affirmatively it suffices to produce a computably
represented structure $\A$ with $\AUT(\A)$ countable such that:
\begin{enumerate}[label=\textup{(\roman*)}]
\item \textup{(P1)} every $\pi\in\AUT(\A)$ is $\Sigma^{in}_1$-definable from
  some finite tuple; and
\item for every finite tuple $\bar a$ there exist $\pi\in\AUT(\A)$ and a
  self-embedding $h$ of $\A$ fixing $\bar a\cup\pi(\bar a)$ pointwise with
  $h\circ\pi\neq\pi\circ h$.
\end{enumerate}
\end{corollary}

\begin{proof}
Item (ii) supplies, via Lemma~\ref{lem:embcrit} with
$\bar c=\bar a\cup\pi(\bar a)$, an automorphism $\pi$ not
$\Sigma^{in}_1$-definable from $\bar a\cup\pi(\bar a)$; this is (P2).  Together
with (i) this is exactly the requirement of Question~1.2.
\end{proof}

\section{Candidate I: the directed/anchored spine ($\AUT\cong\Z$)}\label{sec:candI}

\begin{definition}[The structure $\A_{\mathrm{I}}$]\label{def:AI}
The universe is $A=\{(n,i):n\in\N,\ i\in\Z\}$, in computable bijection with
$\N$.  Put $G_n=\{(n,i):i\in\Z\}$, so $A=\bigsqcup_n G_n$.  Three binary
relations:
\begin{align*}
P\bigl((n,i),(m,j)\bigr)&\iff n=m\ \text{and}\ |i-j|=1,
&&\text{(undirected path in each gadget)}\\
D\bigl((n,i),(m,j)\bigr)&\iff n=m=0\ \text{and}\ j=i+1,
&&\text{(directed anchor on $G_0$)}\\
R\bigl((n,i),(m,j)\bigr)&\iff m=n+1\ \text{and}\ j=i.
&&\text{(directed, position-aligned spine)}
\end{align*}
\end{definition}

\begin{lemma}[$\AUT(\A_{\mathrm{I}})\cong\Z$]\label{lem:AI-aut}
For $k\in\Z$ let $\sigma_k(n,i)=(n,i+k)$.  Then $k\mapsto\sigma_k$ is an
isomorphism $\Z\xrightarrow{\sim}\AUT(\A_{\mathrm{I}})$.  In particular
$\AUT(\A_{\mathrm{I}})$ is countable.
\end{lemma}

\begin{proof}
Each $\sigma_k$ preserves and reflects $P,D,R$ (direct check; e.g.\
$R((n,i),(n+1,i))$ has image $R((n,i+k),(n+1,i+k))$), and
$\sigma_k\sigma_\ell=\sigma_{k+\ell}$ with $\sigma_k=\mathrm{id}\Rightarrow k=0$;
so $k\mapsto\sigma_k$ is an injective homomorphism.  For surjectivity let
$\phi\in\AUT(\A_{\mathrm{I}})$.

\emph{$\phi$ fixes every gadget index.}  In the directed graph $R$ a point has
$R$-in-degree $0$ iff $n=0$, and otherwise $1$; the out-degree is always $1$.
Automorphisms preserve in/out-degrees, so $\phi(G_0)=G_0$; inducting along the
vertex-disjoint $R$-rays $(0,j)\to(1,j)\to\cdots$ gives $\phi(G_n)=G_n$ for all
$n$.

\emph{The shift is uniform.}  Writing $\phi(n,i)=(n,I(n,i))$, applying $\phi$ to
$R((n,i),(n+1,i))$ yields $R((n,I(n,i)),(n+1,I(n+1,i)))$, forcing
$I(n+1,i)=I(n,i)$; hence $I(n,i)=:f(i)$ is independent of $n$, and $f$ is a
bijection of $\Z$.

\emph{$f$ is a translation.}  On $G_1$ (where $P$ is the undirected path),
$|i-(i+1)|=1\Rightarrow|f(i)-f(i+1)|=1$; a bijection of $\Z$ with steps $\pm1$ is
monotone, so $f(i)=i+k$ or $f(i)=c-i$.  On $G_0$ the directed $D$ forces
$f(i+1)=f(i)+1$, excluding $f(i)=c-i$.  Thus $f(i)=i+k$ and $\phi=\sigma_k$.
\end{proof}

\begin{lemma}[$\A_{\mathrm{I}}$ fails (P2): orientation leak]\label{lem:AI-fail}
Every $\sigma_k\in\AUT(\A_{\mathrm{I}})$ is $\Sigma^{in}_1$-definable from the
\emph{empty} tuple.  Hence for every finite $\bar a$ and every $\pi$, $\pi$ is
$\Sigma^{in}_1$-definable from $\bar a\cup\pi(\bar a)$, so (P2) fails.
\end{lemma}

\begin{proof}
It suffices to treat $\sigma=\sigma_1$ (compose for $\sigma_k$, $k\ge1$; reverse
the $D$-chain for $k\le-1$; $k=0$ is ``$x=y$'').  For $n\ge1$ set
\[
\phi_n(x,y):=\exists u_0\cdots u_n\,\exists v_0\cdots v_n\Bigl(
D(u_0,v_0)\wedge\!\!\bigwedge_{t=0}^{n-1}\!\!\bigl(R(u_t,u_{t+1})\wedge R(v_t,v_{t+1})\bigr)
\wedge x=u_n\wedge y=v_n\Bigr),
\]
and $\phi_0(x,y):=D(x,y)$.  Since the $R$-rays are the vertex-disjoint
$(0,j)\to(1,j)\to\cdots$, the only witnesses with $u_n=(n,a)$ are $u_t=(t,a)$,
forcing $u_0=(0,a)$; likewise $v_0=(0,b)$ and $v_n=(n,b)$.  Then $D(u_0,v_0)$
holds iff $b=a+1$, i.e.\ $y=\sigma(x)$ with $x,y\in G_n$.  Hence
$\sigma(x)=y\iff\exists n\,\A_{\mathrm{I}}\models\phi_n(x,y)$, which is
\eqref{star} with empty parameters.  Adding parameters cannot remove
definability, so $\pi$ is definable from any superset, including
$\bar a\cup\pi(\bar a)$.
\end{proof}

\begin{remark}\label{rem:I-diag}
The leak is precisely the failure of Corollary~\ref{cor:spec}(ii): by
Lemma~\ref{lem:embcrit}, $\sigma$ being parameter-free definable means no
self-embedding fixing $\varnothing$ can fail to commute with $\sigma$.  The
directed spine $R$ and anchor $D$ pin a global orientation every self-embedding
must respect.  (An earlier draft ``proved'' (P2) for $\A_{\mathrm{I}}$ by
reflecting a far gadget via a \emph{partial} map defined only on a substructure;
that is invalid, since the witnesses of $\phi_n$ run down the spine to $G_0$ and
the partial map does not act on them.  Lemma~\ref{lem:AI-fail} settles the matter
unconditionally.)
\end{remark}

\section{Candidate II: undirected coupling ($\AUT\cong D_\infty$)}\label{sec:candII}

Deleting the directed anchor and making the coupling undirected removes the
global orientation, replacing $\Z$ by the infinite dihedral group.

\begin{definition}[The structure $\A_{\mathrm{II}}$]\label{def:AII}
The universe is $A=\{(n,i):n\in\N,\ i\in\Z\}$ with gadgets $G_n=\{(n,i):i\in\Z\}$.
To pin the gadget indices we add, for each $n$, a \emph{tag}: a rigid directed
path $t^n_0\to t^n_1\to\cdots\to t^n_{n+1}$ (a fresh copy of the finite linear
order of length $n+2$, with a directed binary relation $E$), together with a
unary predicate $U$ true exactly of the heads $t^n_0$, and a binary
\emph{membership} relation $M$ with $M(t^n_0,(n,i))$ for \emph{all} $i\in\Z$ and
no other instances.  Thus tag $n$ is attached to the \emph{whole} gadget $G_n$
(not to any single point of it), and tags of different lengths are
non-isomorphic.  The internal gadget relations are
\begin{align*}
P\bigl((n,i),(m,j)\bigr)&\iff n=m\ \text{and}\ |i-j|=1,&&\text{(undirected path)}\\
R\bigl((n,i),(m,j)\bigr)&\iff |n-m|=1\ \text{and}\ j=i.&&\text{(undirected, position-aligned)}
\end{align*}
There is no directed relation among the points of the gadgets themselves.  The
structure is computably represented (all relations are decidable on the explicit
universe).
\end{definition}

\begin{lemma}[$\AUT(\A_{\mathrm{II}})\cong D_\infty$]\label{lem:AII-aut}
$\AUT(\A_{\mathrm{II}})=\{\,(n,i)\mapsto(n,\varepsilon i+k):\varepsilon\in\{\pm1\},\,k\in\Z\,\}$,
the infinite dihedral group; it is countable.  (Automorphisms fix every tag
pointwise.)
\end{lemma}

\begin{proof}
A tag is a finite directed path with a marked head $U$ of length $n+2$; the tags
have distinct lengths, so each is rigid and is fixed pointwise by every
automorphism, and the assignment $n\mapsto$ (its tag) is preserved.  The
$M$-relation then forces $\phi(G_n)=G_n$ setwise, because the points $M$-related
to the head $t^n_0$ are exactly $G_n$.  On a single gadget the only structure is
the undirected bi-infinite path $P$, whose automorphism group is
$D_\infty=\{i\mapsto\varepsilon i+k\}$.  Writing the action on $G_n$ as
$i\mapsto\varepsilon_n i+k_n$, preservation of $R((n,i),(n+1,i))$ (undirected,
position-aligned) requires $\varepsilon_n i+k_n=\varepsilon_{n+1}i+k_{n+1}$ for
all $i$, whence $\varepsilon_n=\varepsilon_{n+1}$ and $k_n=k_{n+1}$.  So the sign
and shift are global, giving exactly $D_\infty$.  Conversely each global
$(n,i)\mapsto(n,\varepsilon i+k)$ preserves $P,R,E,U,M$ (it fixes every tag
pointwise), hence is an automorphism.
\end{proof}

\begin{lemma}[(P1) for $\A_{\mathrm{II}}$]\label{lem:AII-P1}
Every $\pi\in\AUT(\A_{\mathrm{II}})$ is $\Sigma^{in}_1$-definable from the two
parameters $a=(0,0)$, $b=(0,1)$.
\end{lemma}

\begin{proof}
The ordered pair $(a,b)$ orients $G_0$: ``$z$ is the point at $P$-distance $i$
from $a$ on the $b$-side'' is, for $i\ge1$, the existential statement
$\exists p_0\cdots p_i\,(p_0=a\wedge p_1=b\wedge\bigwedge_t P(p_t,p_{t+1})\wedge
\bigwedge_{s<t}p_s\neq p_t\wedge p_i=z)$ (a non-backtracking $P$-geodesic), and
$i=0$ is $z=a$.  The position-aligned $R$ transports this to every gadget: ``$z$
is the position-$i$ point of $G_n$'' is an existential $R$-chain of length $n$
from the position-$i$ point of $G_0$.  Hence for each $n$ and each
$(\varepsilon,k)$ there is an existential
$\theta_n^{(\varepsilon,k)}(a,b,x,y)$ defining
$\{((n,i),(n,\varepsilon i+k)):i\in\Z\}$, and $\pi(x)=y\iff\exists
n\,\theta_n^{(\varepsilon,k)}(a,b,x,y)$.
\end{proof}

\begin{lemma}[$\A_{\mathrm{II}}$ fails (P2) at nonempty tuples]\label{lem:AII-fail}
For every \emph{nonempty} finite tuple $\bar a$ and every
$\pi\in\AUT(\A_{\mathrm{II}})$, $\pi$ is $\Sigma^{in}_1$-definable from
$\bar a\cup\pi(\bar a)$.  Hence the negative clause (P2), which requires a
non-definable $\pi$ for \emph{every} finite tuple, fails for $\A_{\mathrm{II}}$.
\end{lemma}

\begin{proof}
Write $\pi:(n,i)\mapsto(n,\varepsilon i+k)$.  If $\pi=\mathrm{id}$ it is defined
by ``$x=y$''; assume $\pi\neq\mathrm{id}$.

\smallskip
\emph{Case $\varepsilon=1$ (translation by $k\neq0$), with $\bar a\neq\varnothing$.}
Pick an entry $a=(n_0,i_0)$ of $\bar a$;
then $a$ and $\pi(a)=(n_0,i_0+k)$ are two points of $G_{n_0}$ at distinct
positions, both in $\bar a\cup\pi(\bar a)$.  As in Lemma~\ref{lem:AII-P1} the
ordered pair $(a,\pi(a))$ orients $G_{n_0}$, and the position-aligned $R$
transports the orientation and the displacement $k$ to every gadget, giving an
existential family for $\pi$ over $\{a,\pi(a)\}\subseteq\bar a\cup\pi(\bar a)$.

\smallskip
\emph{Case $\varepsilon=-1$ (reflection).}  Here $\pi(n,i)=(n,2c-i)$ for the
fixed center position $c:=k/2$ if $k$ is even, and $\pi$ has no fixed points if
$k$ is odd.  We treat both via the \emph{single} parameter structure, using only
that reflections are symmetric, so no orientation is needed.
\begin{itemize}
\item If $k$ is odd, $\pi$ has no fixed point; pick any entry $a=(n_0,i_0)$ of
  $\bar a$, so $\pi(a)=(n_0,2c-i_0)\neq a$.  The pair $\{a,\pi(a)\}$ lies in
  $G_{n_0}$ and its $P$-midpoint (the unique point at equal $P$-distance from
  both on the path) is the half-integer position $c$, i.e.\ there is no center
  vertex but the unordered pair determines $c$; transport by $R$ then defines,
  for every gadget, ``$y$ is the $P$-reflection of $x$ through the midpoint of
  the $R$-aligned copies of $a,\pi(a)$,'' an existential family over
  $\{a,\pi(a)\}$.
\item If $k$ is even, the center $c$ is an integer position.  Choose $a=(n_0,i_0)$
  in $\bar a$.  If $i_0\neq c$ then $\pi(a)\neq a$ and the previous bullet
  applies.  If $i_0=c$ then $a$ \emph{is} a fixed center, and we define $\pi$
  from the single parameter $a$ by
  \[
  \psi_d(a,x,y):=\exists p_0\cdots p_d\,q_0\cdots q_d\Bigl(
  p_0=q_0\wedge \mathrm{ctr}(a,p_0)\wedge\!\bigwedge_t\!\bigl(P(p_t,p_{t+1})\wedge P(q_t,q_{t+1})\bigr)
  \]
  \[
  \wedge\!\bigwedge_{s<t}\!(p_s\neq p_t\wedge q_s\neq q_t)\wedge p_1\neq q_1
  \wedge x=p_d\wedge y=q_d\Bigr),
  \]
  where $\mathrm{ctr}(a,p_0)$ is the existential $R$-chain asserting ``$p_0$ is
  the position-$c$ point of the gadget of $x$,'' i.e.\ the $R$-aligned copy of
  $a$.  For $x$ at $P$-distance $d$ from this center, the two non-backtracking
  geodesics with distinct first steps ($p_1\neq q_1$) end at the two points at
  distance $d$ on opposite sides, namely $x$ and its reflection $y$; hence
  $\bigvee_d\psi_d$ defines $\pi$ over $\{a\}$.
\end{itemize}

In all cases $\pi$ is $\Sigma^{in}_1$-definable from a subset of
$\bar a\cup\pi(\bar a)$, so (P2) fails.
\end{proof}

\begin{remark}[The empty tuple, and why $\A_{\mathrm{II}}$ still fails]\label{rem:empty-gap}
At the \emph{empty} tuple $\A_{\mathrm{II}}$ does admit a (P2)-witness: a
translation $\sigma_1$ is \emph{not} $\Sigma^{in}_1$-definable from $\varnothing$,
because the reflection $\rho:(n,i)\mapsto(n,-i)$ is an automorphism fixing
$\varnothing$ that does not commute with $\sigma_1$, so Lemma~\ref{lem:embcrit}
(with $h=\rho$) applies.  However, Question~1.2 requires a (P2)-witness for
\emph{every} finite tuple, and Lemma~\ref{lem:AII-fail} shows that at every
\emph{nonempty} tuple $\bar a$ \emph{every} $\pi\in\AUT(\A_{\mathrm{II}})$ is
$\Sigma^{in}_1$-definable from $\bar a\cup\pi(\bar a)$.  Hence $\A_{\mathrm{II}}$
fails (P2) and does not answer Question~1.2.  The genuine failure is at nonempty
tuples, where $\bar a\cup\pi(\bar a)$ supplies an orientation (translations) or a
center (reflections).
\end{remark}

\begin{remark}[No useful self-embeddings]\label{rem:II-emb}
A bi-infinite undirected path admits no proper self-embedding: an injective
relation-reflecting homomorphism of $(\Z;|i-j|=1)$ into itself has connected
bi-infinite image, hence is onto, hence an automorphism.  With the
position-aligned $R$ and the rigid tags, every self-embedding of
$\A_{\mathrm{II}}$ is therefore already an automorphism, and pointwise
stabilizers of a center-plus-one-point are trivial.  By
Corollary~\ref{cor:spec} and Remark~\ref{rem:autos} this is exactly why (P2)
fails at nonempty tuples.
\end{remark}

\section{The obstruction, and the target for a positive answer}\label{sec:obstruction}

The two refutations share a mechanism, which we abstract.

\begin{lemma}[Position-aligning couplings leak the orienting datum]\label{lem:obstruction}
Suppose $\A$ is partitioned into gadgets $\{G_n\}_{n\in\N}$, each a bi-infinite
line with a position coordinate $i\in\Z$, such that
\begin{enumerate}[label=\textup{(\alph*)}]
\item every automorphism fixes each gadget index and acts on $G_n$ by an isometry
  $i\mapsto\varepsilon_n i+k_n$ of $\Z$;
\item the cross-gadget relations are \emph{position-aligning}: there is an
  existential family identifying, from a point of one gadget, the
  same-position point of any other gadget; and
\item $\AUT(\A)$ is forced to a single global isometry
  $(\varepsilon,k)$ (equivalently $\AUT(\A)\in\{\Z,D_\infty\}$).
\end{enumerate}
Then for every $\pi\in\AUT(\A)$ and every finite $\bar a$ that meets some gadget
in a way determining the global datum $(\varepsilon,k)$ \textup{(}two distinct
positions of one gadget for a translation; a center, or a symmetric pair, for a
reflection\textup{)}, $\pi$ is $\Sigma^{in}_1$-definable from
$\bar a\cup\pi(\bar a)$.  Consequently no such $\A$ satisfies (P2) for all
nonempty finite tuples.
\end{lemma}

\begin{proof}
By (c) it suffices to recover $(\varepsilon,k)$ and then transport it.  The pair
$\{a,\pi(a)\}\subseteq\bar a\cup\pi(\bar a)$ determines $(\varepsilon,k)$ on the
gadget of $a$: for a translation, two distinct positions give the oriented
displacement $k$; for a reflection, the pair (or a single fixed center) gives
the center, and reflections are symmetric so need no orientation.  By (b) a
length-$n$ existential $R$-chain identifies the same-position point of $G_n$,
yielding for each $n$ an existential $\phi_n(\bar a\cup\pi(\bar a),x,y)$ defining
$\pi\restriction G_n$.  The family $(\phi_n)_n$ is exactly the latitude allowed
by $\Sigma^{in}_1$-definability, establishing \eqref{star}.
\end{proof}

Lemma~\ref{lem:obstruction} pinpoints the difficulty: a coupling strong enough to
force $\AUT(\A)$ \emph{countable} via a single global isometry is automatically
strong enough to \emph{transport} the orienting/value datum along finite
existential chains, defeating (P2).  Escaping this requires, by
Corollary~\ref{cor:spec} and Remark~\ref{rem:autos}, a genuinely different
automorphism group together with proper self-embeddings.

\begin{proposition}[Target specification]\label{prop:target}
Any structure $\A$ answering Question~1.2 affirmatively must have, simultaneously:
\begin{enumerate}[label=\textup{(T\arabic*)}]
\item \emph{(Countable, not a single global isometry.)}  $\AUT(\A)$ is countable
  and, for every finite $\bar a$, contains an automorphism $\pi$ that is the
  identity on a region containing $\bar a$ but nontrivial on a ``free'' region
  disjoint from $\bar a$ --- so that $\pi(\bar a)=\bar a$ and
  $\bar a\cup\pi(\bar a)=\bar a$ carries no information about $\pi$ on the free
  region.  In particular $\AUT(\A)$ is \emph{not} forced to a single global
  isometry; the natural shape is a finite-support group such as
  $\bigoplus_n\Z$ or $\bigoplus_n\Z/L_n$.
\item \emph{(Rich proper self-embeddings.)}  For every finite $\bar a$, the
  self-embeddings fixing $\bar a$ pointwise include a proper one $h$ failing to
  commute with the witnessing $\pi$ of (T1) on the free region
  (Lemma~\ref{lem:embcrit}).
\end{enumerate}
\end{proposition}

\begin{proof}
Immediate from Corollary~\ref{cor:spec}: (T1) is the only way to make
$\bar a\cup\pi(\bar a)$ uninformative while keeping $\AUT(\A)$ countable, since
Lemma~\ref{lem:obstruction} rules out single-global-isometry couplings; and
(T2) is precisely Corollary~\ref{cor:spec}(ii).
\end{proof}

\begin{remark}[The remaining difficulty, stated honestly]\label{rem:open}
Realizing (T1) is the crux.  Uncoupled translatable gadgets give
$\AUT(\A)=\prod_n\Z$ (uncountable); any position-aligning coupling that restores
countability forces a single global isometry (Lemma~\ref{lem:obstruction}) and
destroys (T2).  Forcing \emph{finite support} ($\bigoplus_n$, countable) without
pinning the individual gadgets --- so that cofinitely many gadgets are fixed yet
no gadget is individually rigid --- is achieved by none of the naive devices:
attaching a rigid backbone to a gadget's origin pins that gadget and trivializes
its automorphisms, while a loose attachment fails to bound the support.  A
correct construction must decouple ``cofinitely many gadgets fixed'' from ``each
gadget pinned'' and simultaneously supply the proper self-embeddings of (T2).
We have not produced such a construction, and we do not assert the answer is
\textsc{Yes} on the strength of the Background's framing alone.
\end{remark}

\section{Status}

With complete proofs we have established:
\begin{itemize}
\item the embedding criterion (Lemma~\ref{lem:embcrit}) and the specification
  (Corollary~\ref{cor:spec}) for a positive answer;
\item exact computations $\AUT(\A_{\mathrm{I}})\cong\Z$
  (Lemma~\ref{lem:AI-aut}) and $\AUT(\A_{\mathrm{II}})\cong D_\infty$
  (Lemma~\ref{lem:AII-aut}), both countable;
\item that $\A_{\mathrm{I}}$ satisfies (P1) but fails (P2) at \emph{every} tuple
  (Lemma~\ref{lem:AI-fail}), and that $\A_{\mathrm{II}}$ satisfies (P1) and fails
  (P2) at every \emph{nonempty} tuple (Lemma~\ref{lem:AII-fail},
  Remark~\ref{rem:empty-gap});
\item the structural obstruction (Lemma~\ref{lem:obstruction}) and the precise
  target (Proposition~\ref{prop:target}) a positive solution must realize.
\end{itemize}
What remains \emph{open} is a structure meeting (T1) and (T2) of
Proposition~\ref{prop:target}, i.e.\ a complete affirmative (or negative) answer
to Question~1.2.  The issue is therefore reported as \emph{in progress}, not
resolved.

\end{document}
